\subsection{The space of closed subgroups of G}

Let us denote by \(\Sigma\) the set of closed subgroups of G; if one associates to each measure \(\alpha\in\Gamma\) the subgroup \(H_{\alpha}\) that is the support of \(\alpha\) , one obtains a mapping (called canonical) of \(\Gamma\) into \(\Sigma\) , which is clearly surjective and permits canonically identifying \(\Sigma\) with the set of orbits of the group of homotheties in \(\Gamma\) with ratio \textgreater0. The set \(\Sigma\) , equipped with the quotient topology of the vague topology on \(\Gamma\) , is called the space of closed subgroups of G.

THEOREM 1. --- Let G be a locally compact group. The space \(\Sigma\) of closed subgroups of G is compact. Moreover, one has the following properties:

\begin{enumerate}
\def\labelenumi{(\roman{enumi})}
\item
  The set \(\Sigma^0\) of unimodular closed subgroups of \(\mathbf{G}\) is closed in \(\Sigma\) (hence is compact).
\item
  If G is generated by a compact neighborhood of e, then the set \(\Sigma_{c}^{0}\) of unimodular closed subgroups H of G such that the quotient space G/H is compact, is open in \(\Sigma^{0}\) (hence is locally compact).
\item
  For every relatively compact open neighborhood U of e in G, the set \(D_{U}\) of discrete subgroups H of G such that \(H \cap U = \{e\}\) is closed in \(\Sigma^{0}\) (hence is compact).
\end{enumerate}

It follows from Prop. 3 of No.~1 that \(\Sigma\) is homeomorphic to \(\Gamma_{\varphi}\) , hence is compact by Prop. 2 of No.~1. Moreover, it was noted at the beginning of No.~2 that the set \(\Gamma^0\) of measures \(\alpha \in \Gamma\) such that \(\mathbf{H}_{\alpha}\) is unimodular is closed in \(\Gamma\) ; since \(\Gamma^0\) is stable under the homotheties with ratio \(>0\) , the image \(\Sigma^0\) of \(\Gamma^0\) in \(\Sigma\) is a closed subset of \(\Sigma\) , which proves (i).

Property (ii) will be a consequence of the following proposition:

PROPOSITION 6. --- Suppose that the locally compact group G is generated by a compact neighborhood of e. Then the set \(\Gamma_{c}^{0}\) of measures \(\alpha\in\Gamma^{0}\) such that \(\mathbf{G} / \mathbf{H}_{\alpha}\) is compact is open in \(\Gamma^0\) , and the restriction to \(\Gamma_c^0\) of the mapping \(\alpha \mapsto \| \mu_{\alpha}\|\) is vaguely continuous.

With the notations of Prop. 5 of No.~2, we have, for \(g \in \mathcal{K}_{+}(\mathbf{G})\) ,

\[
\Gamma^ {0} (g) \subset \Gamma_ {c} ^ {0}.\tag{6}
\]

For, if K is the support of g, the relation \(\int g(xs) d\alpha(s) \geqslant 1\) for all \(x \in G\) implies \(KH_{\alpha} = G\) , the integral obviously being zero on the complement of \(KH_{\alpha}\) , therefore \(G/H_{\alpha} = \pi_{\alpha}(K)\) is compact. Given a measure \(\alpha \in \Gamma_{c}^{0}\) , it will therefore suffice to define a function \(g \in \mathcal{K}_{+}(G)\) such that \(\Gamma^{0}(g)\) is a neighborhood of \(\alpha\) in \(\Gamma^{0}\) . Since \(G/H_{\alpha}\) is compact and the canonical mapping \(f \mapsto f_{\alpha}\) of \(\mathcal{K}_{+}(G)\) into \(\mathcal{K}_{+}(G/H_{\alpha})\) is surjective (Ch. VII, §2, No.~2), there exists a function \(g \in \mathcal{K}_{+}(G)\) such that \(\int g(xs) d\alpha(s) = 2\) for every \(x \in G\) . Let K be the (compact) support of g, L a symmetric compact neighborhood of e in G that generates G; the mapping \(\beta \mapsto g*\beta\) of \(\mathcal{M}_{+}(G)\) into \(\mathcal{C}(G)\) being vaguely continuous (§4, No.~2, Remark 1), there exists a neighborhood W of \(\alpha\) in \(\Gamma^{0}\) such that

\[
(g * \beta) (x) = \int g (x s) d \beta (s) \geqslant 1\tag{7}
\]

for all \(\beta \in W\) and \(x \in LK\) . The first member of (7) being equal to zero outside \(\mathrm{KH}_{\beta}\) , the relation \(\beta \in W\) implies

\[
\mathbf {L K} \subset \mathbf {K H} _ {\beta},
\]

from which one deduces, by induction on n, that \(L^{n}K \subset KH_{\beta}\) for every integer n \textgreater{} 0; since L generates G, we therefore have \(G = KH_{\beta}\) for every measure \(\beta \in W\) , which proves that \(W \subset \Gamma_{c}^{0}\) . On the other hand, the first member of (7) being invariant on the right under \(H_{\beta}\) , the inequality (7) is also valid for \(x \in LKH_{\beta} = G\) ; therefore \(W \subset \Gamma^{0}(g)\) , which proves the proposition.

Finally, (iii) will be a consequence of the following proposition:

PROPOSITION 7. --- Let \(N \subset \Gamma^{0}\) be the subspace of normalized Haar measures on the discrete subgroups of G, and for every relatively compact open neighborhood U of e in G, let \(N_{U}\) be the subset of N formed by the \(\alpha\) such that \(H_{\alpha} \cap U = \{e\}\) . Then:

\begin{enumerate}
\def\labelenumi{\alph{enumi})}
\item
  \(\mathbf{N}_{\mathrm{U}}\) is compact.
\item
  The interiors of the sets \(\mathbf{N}_{\mathbf{U}}\) in \(\mathbf{N}\) form a covering of \(\mathbf{N}\) , as \(\mathbf{U}\) runs over the set of relatively compact open neighborhoods of \(e\) in \(\mathbf{G}\) .
\item
  For a subset M of N to be relatively compact in N, it is necessary and sufficient that there exist a relatively compact open neighborhood U of e in G such that \(M \subset N_{U}\) .
\end{enumerate}

Since \(D_{U}\) is the image of \(N_{U}\) under the canonical continuous mapping \(\Gamma \rightarrow \Sigma\) , the assertion (iii) of Th. 1 will result at once from Prop. 7 a).

To prove Prop. 7, we observe that \(N_{U}\) can be defined as the subset of \(\Gamma^{0}\) formed by the \(\alpha\) such that both

\[
\alpha (\{e \}) \geqslant 1 \quad \text { and } \quad \alpha (U) \leqslant 1.
\]

Now, if A is compact (resp. open and relatively compact) in G, then the mapping \(\alpha \mapsto \alpha(A)\) of \(\mathcal{M}_{+}(G)\) into R is upper (resp. lower) semicontinuous for the vague topology (Ch. IV, §4, No.~4, Cor. 3 of Prop. 5 and loc. cit., §1, No.~1, Prop. 4); we thus see that \(N_{U}\) is a closed subset of \(\Gamma^{0}\) . Moreover, let \(\varphi \in \mathcal{K}_{+}(G)\) be a function such that \(\varphi(e) = 1\) and \(\varphi(x) = 0\) on G - U; it is clear that \(\int \varphi(x)d\alpha(x) = 1\) for all \(\alpha \in N_{U}\) ; Prop. 2 of No.~1 therefore shows that \(N_{U}\) is a compact set, which proves a). On the other hand let V be a relatively compact open neighborhood of \(e\) in G such that \(\overline{V} \subset U\) , and let \(\varphi \in \mathcal{K}_{+}(G)\) , with support contained in U and such that \(\varphi(x) = 1\) on V. Then \(\alpha(\varphi) = 1\) for \(\alpha \in N_{U}\) , therefore there exists a neighborhood W of \(\alpha\) in N such that \(\beta(\varphi) < 2\) for \(\beta \in W\) ; it is then clear that \(W \subset N_{V}\) , therefore \(N_{V}\) is a neighborhood of \(N_{U}\) . Since the \(N_{U}\) cover N, this proves b). Finally, every compact subset M of N is contained in a finite union of sets \(N_{U_i}\) ( \(1 \leqslant i \leqslant n\) ), and since \(\bigcup_{i} N_{U_i} \subset N_{U}\) , where \(U = \bigcap_{i} U_i\) , this proves c).

COROLLARY. --- The subspace N of \(\Gamma^0\) is locally compact.
% semantic-fragments: legacy-input-seam
\relax{}
